\chapter*{Resumen}
\addcontentsline{toc}{chapter}{Resumen}

La seguridad alimentaria global depende de la estabilidad productiva de cultivos básicos como maíz, arroz, trigo y soya. Sin embargo, la variabilidad climática asociada a modos de gran escala, como El Niño-Oscilación del Sur (ENSO) y la Oscilación Madden-Julian (MJO), puede alterar la disponibilidad hídrica, la temperatura y la ocurrencia de eventos extremos durante fases críticas del desarrollo agrícola. Esta tesis propone caracterizar la relación entre la variabilidad climática global y el rendimiento de cereales, con el fin de identificar regímenes agroclimáticos y umbrales climáticos aplicables al diseño de seguros paramétricos agrícolas.

El documento se estructura en dos artículos científicos. El primero analiza la sensibilidad espacio-temporal del rendimiento de maíz, arroz, trigo y soya frente a ENSO y MJO, integrando anomalías climáticas, rendimientos agrícolas y calendarios fenológicos globales. El segundo utiliza los insumos derivados del primer artículo para proponer una metodología de identificación de umbrales climáticos y diseño preliminar de disparadores paramétricos orientados a la gestión del riesgo agrícola y al fortalecimiento de la seguridad alimentaria.

% TODO: completar con metodología definitiva, fuentes de datos, periodo de análisis y principales resultados propios.

\textbf{Palabras clave:} seguridad alimentaria; ENSO; MJO; cereales; rendimiento agrícola; seguros paramétricos; riesgo climático; fenología agrícola.
